%% kapa: exam=1 round=25 year=2014 period=1 form=A subject=economics subject_ko=경제원론 seq=16 number=56
%% source: Q-Net 감정평가사 기출문제 (cst00309) artlSeq=5101207 file=extracted/1008/5101207/2014년 제25회 감정평가사 1차시험 기출문제/1교시(A형).hwp
%% conversion: kordoc 4.24.51 -> tools/gen1.py (auto); review=proofread-sonnet; flags=none
\begin{question}{56}
지리적으로 분리되어 시장간 전매가 불가능한 두 시장 A, B에서 판매하고 있는 독점기업에 대한 수요곡선이 각각 $P_{ A } =-Q_{ A } + 20$이고, $P_{ B } =-0.5Q_{ B } +10$이다. 한계비용이 5이고 이윤극대화를 추구하는 이 기업의 두 시장에서의 가격은 각각 얼마인가?
\begin{choices}
  \item $P_{ A }$=8, $P_{ B }$=12
  \item $P_{ A }$=12.5, $P_{ B }$=7.5
  \item $P_{ A }$=12, $P_{ B }$=8
  \item $P_{ A }$=7.5, $P_{ B }$=12.5
  \item $P_{ A }$=14, $P_{ B }$=6
\end{choices}
\end{question}
\answer{2}
